%% =====================================================================
%%  CHAPTER 2: PROBLEM AND OBJECTIVES               (budget 4--5 pages)
%%  chapters/ch2-problem.tex
%% =====================================================================
\pcChapter{problem}{5}{הגדרת הבעיה והיעדים}
\begin{pcguide}[מה כותבים כאן]
פרק הבעיה מאפיין את הבעיה \textbf{ללא פתרון} (\latin{5W1H}), מנסח שאלת מחקר והשערות עבודה, מפרק את הסוגיה לעץ
\latin{MECE}, ומסכם הכול במפרט דרישות כמותי: לכל דרישה ערך יעד ושיטת בדיקה. הפתרון עצמו
מופיע רק בפרק~\ref{ch:solution}.
\end{pcguide}

\section{הגדרת הבעיה בשיטת \latin{5W1H}}\label{sec:5w1h}
הבעיה היא אובדן היכולת להפוך מסמכים וניסיון אנושי לידע מקצועי שאפשר לאתר, ללמד ולבחון בזמן קצר, ובמיוחד הקושי
להעביר שיקול דעת מרחבי, סדר פעולות ומקרים חריגים שאי אפשר לבטא היטב בטקסט בלבד. אין לבלבל בין הבעיה לבין היעדר
כלי תוכנה: גם מאגר דיגיטלי עשוי להכיל מידע רב בלי שאפשר להסיק ממנו הדרכה נכונה.

\textbf{הקשר מאקרו.} התעשייה האווירית לישראל \latin{(IAI)} מעסיקה לפי המידע הציבורי כ-\latin{15,000} עובדים בארבע
חטיבות \cite{iai2025about}. בארגון בהיקף כזה כל תהליך חפיפה חוזר על עצמו בהרבה יחידות ותפקידים, ולכן גם שיפור מתון
בהכנת חומרי ההדרכה הוא בעל פוטנציאל השפעה רחב. נתון זה מתאר את סביבת הבעיה בלבד: הוא אינו מצביע על שיעור
עזיבה, על משך חפיפה או על היקף הידע שאבד, ולא התקבלו נתונים כאלה מן הארגון. טבלה~\ref{tab:5w1h} מפרקת את הבעיה
לשש שאלות ומצמידה לכל אחת מדד ביצוע מרכזי \latin{(KPI)} שייאסף בפיילוט.
\begin{pcguide}[כלל]
בשלב זה מתארים את הבעיה בלבד. מציעים פתרון --- רק בפרק הפתרון. הגדירו גם \textbf{מדד} ומצב נוכחי:
למשל “שיעור אספקה בזמן \pct{70}, והיעד \pct{85}”. אל תכתבו “חסרה מערכת כלשהי”: הגדירו את צוואר הבקבוק ולא את הפתרון.
\end{pcguide}
\begin{table}[!htbp]
\centering
\caption{הגדרת הבעיה בשיטת \latin{5W1H}; המצב הארגוני יימדד בפיילוט}\label{tab:5w1h}
\begin{pctable}{H{0.17},H{0.46},H{0.37}}
  \pchead{שאלה & תשובה (ללא פתרון) & מדד איסוף מוצע}
  \pcrow{\textbf{מה} \latin{(What)} & פער בין מסמכים, לקחי מומחים ויכולת התרגול וההחלטה של עובד חדש & זמן הכנת יחידת הדרכה; שיעור ידע חסר; נכונות פעולה בתרחיש}
  \pcrow{\textbf{איפה} \latin{(Where)} & במעבר בין יצירת ידע, פרישה, חפיפה ובדיקת כשירות ביחידות ההדרכה והרכש של \latin{IAI} & מפת תהליך לפי יחידה ותפקיד}
  \pcrow{\textbf{מתי} \latin{(When)} & במעבר תפקיד, בשינוי גרסה של נוהל, בתקלה או בתרחיש עומס & תדירות עדכון מסמך; זמן המתנה לאישור}
  \pcrow{\textbf{מי} \latin{(Who)} & מקבל ההחלטות בארגון \latin{(Decision Maker)}: ראש צוות מערכות מידע תפעוליות ב-\latin{IAI}, איש הקשר לפרויקט. משתמשים: עובדים חדשים ומדריכים. בעלי עניין: מומחי תוכן \latin{(SME)} ואבטחת מידע & מספר משתתפים ותפקידים בפיילוט; מועד אישור הגישה}
  \pcrow{\textbf{למה} \latin{(Why)} & עלות זמן מומחים, אי-אחידות בהדרכה וסיכון לטעויות & שעות עבודה, שגיאות ומידת ביטחון}
  \pcrow{\textbf{איך} \latin{(How)} & תיעוד, חפיפה ידנית וקובצי הדרכה במצב הבסיס & מדידת התהליך הידני לפני הפיתוח}
\end{pctable}
\end{table}
זהות מקבל ההחלטות נגזרת מתפקידו של איש הקשר של \latin{IAI} בפרויקט, והיא תאושר מולו בפגישת ההיכרות (פרק~\ref{ch:summary}).
הנתונים ייאספו רק לאחר הסכמה ואישור מתאים. יחידת המדידה בפיילוט היא הכנת יחידת לימוד אחת מתוך קורפוס לא מסווג,
ולא מספר העובדים בחברה כולה.

\section{שאלת מחקר והשערות עבודה}\label{sec:rq}
\begin{pcguide}[נוסחת שאלת המחקר ובדיקת \latin{SMART}]
“כיצד ניתן \ph{לשפר} את \ph{מדד} מ-\ph{ערך נוכחי} ל-\ph{ערך יעד} באמצעות \ph{שיטה}, בהשוואה ל-\ph{בסיס}?”
בדקו כל שאלה מול \latin{SMART} (טבלה~\ref{tab:smart}) וגם מול \latin{FINER}.
בפרויקט שמטרתו לתכנן מערכת ולא לבחון קשר בין גורמים --- ההשערות מנוסחות כיעדים מדידים.
\end{pcguide}
שאלת המחקר מתמקדת במה שניתן למדוד בפיילוט בן \latin{14} שבועות על קורפוס לא מסווג, ולא בתוצאה ארגונית רחבה
שאין עליה נתונים.
\begin{pcbox}[pcPrimary]{שאלת המחקר}
האם צינור \latin{RAG} מקומי, המפיק יחידות הדרכה ותצורות תרחיש \latin{3D} מחומרים לא מסווגים ומראיונות מומחים ובאישור
מומחה, יכול --- בפיילוט בן \latin{14} שבועות ובהשוואה להפקה ידנית --- (א)~להעלות את דיוק חילוץ הידע מראיונות
מומחים, בבדיקה מול מקורות מאומתים, מ-\pct{60} ל-\pct{90}; (ב)~לקצר בלפחות \pct{50} את זמן הפקתה של יחידת הדרכה
מאושרת; ו-(ג)~להשיג נאמנות עובדתית של לפחות \pct{90} ושיעור השלמת משימות של לפחות \pct{80} בתרחיש ה\threeDadj{}?
\end{pcbox}
\textbf{הגדרות המדדים.} \textbf{דיוק חילוץ הידע}: שיעור הטענות שחולצו מראיונות מומחים ואושרו מול מקור תקף בידי שני מעריכים
\mth{(A_{ext}=N_{verified}/N_{extracted})}, בהתאם ל-\req{16}. \textbf{זמן הפקה}: חציון הדקות מקליטת פריט המקור ועד ליחידה מאושרת,
בהשוואה למשימה שקולה בהכנה ידנית (\req{04}). \textbf{נאמנות עובדתית}: שיעור הטענות ביחידות המאושרות הנתמכות במקור
שצוטט ובגרסתו התקפה (\req{03}). \textbf{השלמת משימות}: שיעור משימות התרחיש שהושלמו ללא עזרת מנחה (\req{15}).

\textbf{מעמד המספרים.} ערך הבסיס \pct{60} לדיוק החילוץ הוא הנחת תכנון של הצוות ולא תוצאה שנמדדה ב-\latin{IAI};
הוא יימדד בשבועות \rng{1}{4} (\latin{D4}), ואם יתברר שונה --- יעד השיפור יוגדר מחדש באותו יחס. הערכים \pct{90},
\pct{50} ו-\pct{80} הם יעדי תכן. הפיילוט אינו בוחן אף טענה על התנהגות הארגון כולו.

\textbf{שאיפה ארגונית ארוכת טווח (הנחה, מחוץ למדידת הפיילוט).} בתכנון הראשוני העריך הצוות כי זמן החפיפה למשימת יעד אחת,
כ-\latin{14} ימי עבודה, עשוי להתקצר לכ-\latin{4} ימים (\numto{14}{4}) אם שלושת רכיביו יקוצרו: המתנה לחומר הדרכה מאושר
\mth{(T_{prep})}, לימוד ומבחן \mth{(T_{learn})} ותרגול מודרך \mth{(T_{practice})}. זמן חפיפה כזה לא נמדד בתעשייה האווירית,
ולכן אינו חלק מהשאלה, מההשערות או ממדדי ההצלחה. הפיילוט יבדוק רק את הרכיבים שהצוות שולט בהם: \mth{T_{prep}} (השערה \hyp{2})
ואיכות הלמידה (\hyp{3} ו-\hyp{4}); אם יוכח שהם עומדים ביעדים, אפשר יהיה לבחון בעתיד, בהסכמה ובנתוני אמת, גם את זמן החפיפה.
טבלה~\ref{tab:smart} בודקת את השאלה מול \latin{SMART}, וטבלה~\ref{tab:hyp} מפרטת את ההשערות.
\begin{table}[!htbp]
\centering
\pcCompactTables
\caption{בדיקת \latin{SMART} לשאלת המחקר}\label{tab:smart}
\begin{pctable}{c,H{0.38},H{0.42}}
  \pchead{קריטריון & שאלת בדיקה & תשובה}
  \pcrow{\latin{Specific} & על מה בדיוק השאלה? & חילוץ ידע מראיונות מומחים והפקת יחידת הדרכה מאושרת ותרחיש \latin{3D} מתוך מסמכים פומביים}
  \pcrow{\latin{Measurable} & איך נדע שהצלחנו (מדד, ערך נוכחי, ערך יעד)? & דיוק חילוץ \pctto{60}{90}; זמן הפקה: חיסכון \gesym{50\%}; נאמנות \gesym{90\%}; השלמת משימות \gesym{80\%}}
  \pcrow{\latin{Achievable} & האם זה ריאלי במשאבים שלנו? & אב-טיפוס מוגבל לתפקיד, לקורפוס ולתרחיש מוגדרים; מודל מקומי וסצנה גנרית אחת}
  \pcrow{\latin{Relevant} & למי הבעיה כואבת? & יחידת ההדרכה, מומחים פורשים ועובדים חדשים: צוואר בקבוק בחפיפה ואובדן ידע מקצועי}
  \pcrow{\latin{Time-bound} & עד מתי, בתוך שנת הסמינריון? & אב-טיפוס וניסוי בתוך \latin{14} שבועות של סמסטר ב' (פרק~\ref{ch:plan})}
\end{pctable}
\end{table}
\begin{table}[!htbp]
\centering
\pcCompactTables
\caption{השערות עבודה: נגזרות מהספרות, מדידות וניתנות להפרכה}\label{tab:hyp}
\begin{pctable}{C{0.07},H{0.31},C{0.12},H{0.18},H{0.32}}
  \pchead{\# & ניסוח ההשערה (גורם, כיוון, מדד) & מדד & מקור בספרות & בדיקה והפרכה}
  \pcrow{\hyp{1} & דיוק חילוץ הידע מראיונות מומחים יעלה מבסיס מוערך של \pct{60} לדיוק של \pct{90} לפחות & שיעור טענות מאושרות & \cite{zuin2025tacit,finkel2026unlocking} & שני מעריכים בודקים מדגם טענות מול מקור תקף; הפרכה אם הדיוק נמוך מ-\pct{90}}
  \pcrow{\hyp{2} & זמן ההפקה החציוני של יחידה מאושרת יהיה קטן או שווה למחצית זמן הבסיס הידני & יחס זמנים & פער: לא נמדד בספרות שנסקרה \cite{hua2025civask} & משימות שקולות עם טופסי מדידת זמן; הפרכה אם היחס גבוה מ-\num{0.5}}
  \pcrow{\hyp{3} & לפחות \pct{90} מהטענות העובדתיות ביחידות המאושרות ייתמכו במקור תקף ובגרסה הנכונה, ולפחות \pct{80} משאלות ההערכה יעמדו במחוון האיכות & נאמנות; שיעור אישור & \cite{es2024ragas,saad2024ares,ramkumar2026pedagogical} & שני מעריכים וטבלת מחוון בן ארבעה ממדים; הפרכה אם אחד מהשיעורים נמוך מהסף}
  \pcrow{\hyp{4} & לפחות \pct{80} ממשימות התרחיש ה\threeDadj{} יושלמו ללא עזרה, ובהשוואה לגרסה \mbox{דו-ממדית} זהה בתוכן יימדד הפרש של \gesym{10} נקודות אחוז & שיעור השלמה; הפרש נק' אחוז & \cite{bodding2025mr}; השוואה זהת-תוכן חסרה (סעיף~\ref{sec:gap}) & ההשלמה: מדד קבלה; ההשוואה: גישושית. כישלון או פגיעה בנגישות יובילו לשינוי ממשק ולא לשינוי נוהל}
\end{pctable}
\end{table}
ההשערות הן יעדי תכן ולא ממצאים. שאלת המחקר מתפרקת אליהן כך: חלק~(א) --- \hyp{1}; חלק~(ב) --- \hyp{2};
חלק~(ג) --- \hyp{3} ו-\hyp{4}. \hyp{4} היא השערת משנה גישושית ואינה מחליפה מבחן כשירות רשמי.

\section{עץ סוגיות}\label{sec:issuetree}
לצורך פירוק הבעיה נבחרו ארבעה ענפים לפי שלבי שרשרת הערך: קליטת הידע ומקורותיו, ניהול גרסאות ועדכוני ידע, הפקת
יחידת ההדרכה, ולמידה והוכחת כשירות (איור~\ref{fig:issuetree}). כל עלה מציין את הנתון או המדד שישיב על השאלה.
\begin{figure}[!htbp]
\centering
\pcFig{issue_tree}{%
\pcAutoFit{\begin{pctikz}[x=1cm,y=1cm]
  \tikzset{lf/.style={pcleaf,text width=4.3cm,font=\footnotesize\sffamily,inner sep=3pt},
           br/.style={pcnode,text width=3.1cm,font=\footnotesize\sffamily\bfseries}}
  \node[pcnodeD,text width=2.5cm] (q) at (0,0) {\tn{פער בחפיפה ובשימור ידע}};
  % branch 1: capture
  \node[br] (a1) at (-4.5,4.8) {\tn{א. קליטת ידע ומקורות}};
  \node[lf] (b1x) at (-10.6,5.6) {\tn{אילו טענות חולצו מהמומחים? \textbf{מדד:} דיוק חילוץ מול מקור תקף, \pctto{60}{90}}};
  \node[lf] (b1y) at (-10.6,4.0) {\tn{מה מותר לעיבוד? \textbf{נתון:} סיווג והרשאה לכל מקור, \pct{100} מתועד}};
  % branch 2: version control
  \node[br,fill=pcWarn!15,draw=pcWarn] (a2) at (-4.5,1.6) {\tn{ב. ניהול גרסאות ועדכוני ידע}};
  \node[lf] (b2x) at (-10.6,2.4) {\tn{איזו גרסה תקפה? \textbf{נתון:} שיוך מקור, עמוד וגרסה לכל טענה, \pct{100}}};
  \node[lf] (b2y) at (-10.6,0.8) {\tn{מה מושפע מעדכון נוהל? \textbf{מדד:} זמן מעדכון ועד ביטול היחידות התלויות, \lesym{1} יום עבודה}};
  % branch 3: generation
  \node[br] (a3) at (-4.5,-1.6) {\tn{ג. הפקת יחידת הדרכה}};
  \node[lf] (b3x) at (-10.6,-0.8) {\tn{האם הפלט מעוגן במקור? \textbf{מדד:} נאמנות עובדתית \gesym{90\%}; הימנעות נכונה \gesym{90\%}}};
  \node[lf] (b3y) at (-10.6,-2.4) {\tn{כמה זמן עד אישור? \textbf{מדד:} חציון דקות מול בסיס ידני, חיסכון \gesym{50\%}}};
  % branch 4: learning
  \node[br] (a4) at (-4.5,-4.8) {\tn{ד. למידה והוכחת כשירות}};
  \node[lf] (b4x) at (-10.6,-4.0) {\tn{האם הלומד משלים את המשימות? \textbf{מדד:} שיעור השלמה ללא עזרה, \gesym{80\%}}};
  \node[lf] (b4y) at (-10.6,-5.6) {\tn{מה מוסיף התרגול ה\threeDadj{}? \textbf{מדד:} הפרש נק' אחוז מול גרסה \mbox{דו-ממדית}, \gesym{10}}};
  \foreach \i in {1,2,3,4}{
    \draw[pcflow] (q.west) -- ++(-0.7,0) |- (a\i.east);
    \draw[pcflow] (a\i.west) -- ++(-0.6,0) |- (b\i x.east);
    \draw[pcflow] (a\i.west) -- ++(-0.6,0) |- (b\i y.east);}
\end{pctikz}}}
\caption{עץ סוגיות \latin{MECE}: ארבעה ענפים לפי שלבי שרשרת הערך; ענף ב' (גרסאות ועדכוני ידע) הוא דרישת הליבה של \latin{IAI}; כל עלה מפנה לנתון או למדד}\label{fig:issuetree}
\end{figure}
\begin{pcguide}[כללי בניית העץ]
כותבים את השאלה הראשית בקצה הימני; מחלקים לשניים עד ארבעה ענפי \latin{MECE}; מפרקים כל ענף עד לשאלה
שאפשר לענות עליה בנתונים; בודקים כל ענף ושואלים מה הנתון שחסר לנו.
\end{pcguide}
\textbf{בדיקת \latin{MECE} וניתוח העץ.} החלוקה היא לפי שלבי הזמן של הידע בארגון: ידע נקלט (א), נשמר ומתעדכן (ב), הופך ליחידת הדרכה
(ג) ומגיע ללומד (ד). כל שאלה שייכת לענף אחד בלבד: שאלת ה”גרסה” רק בענף ב’, איכות התוכן שמופק (נאמנות, זמן הפקה) רק בענף ג’,
ותוצאי הלומד (השלמה, תרומת ה-\latin{3D}) רק בענף ד’, ולכן אין חפיפה בין הפקת התוכן להערכת הלמידה. ענף ב’ מוצג בנפרד משום ששינוי גרסה גורם
ליחידות קיימות להפוך ללא תקפות, והוא לב דרישת \latin{IAI}: עדכון אוטומטי של חומרי ההדרכה לאחר שינוי נוהל. הענף המכריע לפיילוט הוא ג’,
שבו נבחנות \hyp{2} ו-\hyp{3}; ענפים א’ ו-ב’ הם תנאי סף לו (ללא מקור מאושר ובגרסה תקפה אין על מה להפיק), ענף ד’ בוחן את \hyp{4}, ו-\hyp{1} נבחנת בענף א’.
שגיאה בשלב מוקדם משפיעה על המאוחרים, ולכן קשרי התלות נרשמים בניתוח ולא מטושטשים בחלוקה.

\section{מפרט דרישות והאילוצים}\label{sec:requirements}
הדרישות בטבלה~\ref{tab:req} הן חוזה התכן של אב-הטיפוס, וייבדקו על יחידות מידע פומביות: דרישות \textbf{תפקודיות} \latin{(Functional)} המתארות מה המערכת עושה,
ודרישות \textbf{לא-תפקודיות ואילוצים} \latin{(Non-Functional / Constraints)} המתארות כמה טוב ובאילו גבולות. שיטות האימות: בדיקה \latin{(T)},
ניתוח \latin{(A)}, הדגמה \latin{(D)} ובחינה \latin{(I)}.
\begin{pcguide}[דרישות הנדסיות]
כל דרישה: מזהה ייחודי, ניסוח כמותי, סיווג, ערך יעד ושיטת אימות. שיטות אימות:
בדיקה \latin{(T)}, ניתוח \latin{(A)}, הדגמה \latin{(D)}, בחינה \latin{(I)}. הדרישות הן החוזה מול המרצה,
ולכן הפתרון המוצע (פרק~\ref{ch:solution}) ותוכנית הבדיקה שלו מתייחסים לכל מזהה.
בדוח מלא: \rng{8}{12} דרישות, ולפחות אחת מכל סיווג; אילוצי אבטחה וסביבת פריסה נכנסים כאן כאילוצים.
\end{pcguide}
\begin{table}[!htbp]
\centering
\pcCompactTables
\caption{מפרט דרישות הנדסי (\latin{16} דרישות): תפקודיות מול לא-תפקודיות ואילוצים}\label{tab:req}
\begin{pctable}{c,H{0.41},C{0.11},H{0.24},C{0.08}}
  \pchead{מזהה & דרישה & סוג משנה & ערך יעד & אימות}
  \midrule\multicolumn{5}{c}{\texthebrew{\textbf{א. דרישות תפקודיות} \latin{(Functional)}}}\\\midrule
  \pcrow{\req{01} & קליטת מסמכי \latin{PDF} וגרסאות למאגר המקומי & תפקודית & \pct{100} מקובצי הניסוי & \latin{T}}
  \pcrow{\req{02} & שיוך מקור, עמוד וגרסה לכל טענה & תפקודית & \pct{100} מהפריטים & \latin{I}}
  \pcrow{\req{05} & זיהוי שאלות שאין להן מקור והימנעות ממענה & תפקודית & לפחות \pct{90} במדגם & \latin{T}}
  \pcrow{\req{07} & הפקת תרגילי למידה ושאלות לכל יחידה & תפקודית & לפחות \latin{3} שאלות ליחידה & \latin{I}}
  \pcrow{\req{11} & סימון וביטול פרסום של יחידות תלויות לאחר שינוי גרסת מקור & תפקודית & עד יום עבודה אחד & \latin{T}}
  \pcrow{\req{13} & הפקת תצורת \latin{JSON} של תרחיש מיחידה מאושרת, והזנתה לתבנית \latin{3D} פרמטרית & תפקודית & סצנה אחת עם \latin{3} נקודות החלטה & \latin{D}}
  \pcrow{\req{15} & זיהוי אוטומטי של השלמת כל משימת תרחיש לפי תנאי הצלחה בתצורה, ודיווח שיעור השלמה & תפקודית & תנאי הצלחה לכל משימה; סף קבלה \gesym{80\%} השלמה & \latin{T}}
  \midrule\multicolumn{5}{c}{\texthebrew{\textbf{ב. דרישות לא-תפקודיות ואילוצים} \latin{(Non-Functional / Constraints)}}}\\\midrule
  \pcrow{\req{03} & נאמנות עובדתית של יחידות מאושרות & ביצועים & לפחות \pct{90} & \latin{T}}
  \pcrow{\req{04} & זמן הפקת יחידה לעומת הכנה ידנית & ביצועים & חיסכון של \pct{50} לפחות & \latin{A}}
  \pcrow{\req{06} & ממשק אישור: מומחה מאשר או דוחה כל פריט לפני פרסום & ממשק & כל הפריטים & \latin{D}}
  \pcrow{\req{08} & שליפת מקור לשאילתה רגילה & ביצועים & חציון עד \latin{10} שניות & \latin{T}}
  \pcrow{\req{09} & עיבוד מקומי, ללא שליחת תוכן לענן & אילוץ & אפס בקשות לרשת חוץ & \latin{T}}
  \pcrow{\req{10} & בקרת הרשאות ותיעוד פעולות & אילוץ & \pct{100} מפעולות העריכה & \latin{I}}
  \pcrow{\req{12} & איכות פדגוגית של שאלות מאושרות & ביצועים & לפחות \pct{80} & \latin{I}}
  \pcrow{\req{14} & ממשק רישום: החלטות, זמן ומשוב בסימולציה & ממשק & כל פעולות המדגם & \latin{T}}
  \pcrow{\req{16} & דיוק חילוץ הידע מראיונות מומחים & ביצועים & \pctto{60}{90} (בסיס מוערך) & \latin{T}}
\end{pctable}
\end{table}
\req{16} נוספה כדי למדוד ישירות את רכיב שימור הידע בשאלת המחקר. \req{15} מנוסחת כיכולת של המערכת: הסימולציה
חייבת לדעת לקבוע, לפי תנאי הצלחה מוגדרים, אם משימה הושלמה; שיעור ההשלמה של המשתתפים (\hyp{4}) הוא ממצא הפיילוט
המחושב בעזרת יכולת זו, ולא דרישה מן הלומד.

\subsection*{אילוצים}
האילוצים בטבלה~\ref{tab:constraints} קובעים מה אפשר לבנות בכלל, ומתורגמים לדרישות \req{09} ו-\req{10} ולבטיחות תרגול העומס המדומה.
\begin{table}[!htbp]
\centering
\pcCompactTables
\caption{אילוצי הפרויקט לפי קטגוריה}\label{tab:constraints}
\begin{pctable}{H{0.26},H{0.40},H{0.34}}
  \pchead{קטגוריית אילוץ & תיאור וערך & השפעה על התכן}
  \pcrow{אבטחת מידע וסיווג & אין גישה לחומר מסווג; חומר ארגוני רק לאחר אישור וסיווג & פיילוט על מידע פומבי בלבד}
  \pcrow{סביבת פריסה (מקומית, מנותקת, ענן) & אסור לשלוח מידע ארגוני לשירותים חיצוניים & רכיבי אחזור ומודל מקומיים; רשת יוצאת חסומה בניסוי}
  \pcrow{זמינות הנתונים & לא התקבל קורפוס פנימי & שימוש במדריך \latin{FAA} כחלופה לא מסווגת}
  \pcrow{זמן ותקציב & שני חברי צוות, \latin{14} שבועות & היקף ניסוי מצומצם; סצנה גנרית אחת}
  \pcrow{אילוצים ארגוניים ורגולטוריים & בטיחות ופרטיות המשתתפים & תרגול מפתיע רק בהסכמה ובתיאום; ללא אזעקת אמת, ללא הזעקה פיזית וללא מידע אישי מזהה}
\end{pctable}
\end{table}
\pcChapterEnd{problem}

\pcIfSolution{afterproblem}
